\chapter{工作记忆认知总线：从局部暂存到跨尺度显式认知界面}

\begin{chapteroverviewbox}
在前一章中，我们已经由工作记忆有限容量与有限有效认知复杂度这一基本约束得到一个关键结论：一个持续智能体不可能把自身全部内部状态、全部长期知识、全部身体状态、全部环境信息以及全部未来可能性同时展开在当前认知空间中。由此，工作记忆 $h(t)$ 的理论地位必须发生改变。它不能继续被理解为一种容量较小的“短期存储区”，也不能等同于大语言模型中的上下文窗口；更准确地说，$h(t)$ 是多时间尺度认知结构在当前时刻发生\textbf{显式交汇、竞争、组合与行动闭合}的稀疏认知界面，是一个持续智能体内部的\textbf{认知总线（Cognitive Bus）}。

本章将在不依赖后续具体调度模块实现的前提下，对这一命题进行严格展开。我们将从 $h(t)$ 的本体地位出发，说明世界状态、情景记忆 $E$、结构记忆 $\Theta$、自我结构 $\Psi$、外部认知载体 $Me$ 以及行动系统为什么都必须通过一种有限的显式工作面发生跨尺度耦合；进一步说明“进入 $h$”并不意味着复制全部信息，而意味着某个 Agent-CSO 获得当前时刻的全局显式处理权。由此可以得到一个贯穿 AgentOS 原理论的重要认识：\textbf{成熟智能的目标并不是让越来越多内容进入工作记忆，而是让越来越少、越来越关键的结构在正确时刻进入工作记忆。}

本章因此承担一个理论转折功能：此前的三相记忆理论回答“认知结构怎样跨时间存在”，工作记忆有限容量定理回答“当前认知为什么必然稀疏”，而本章进一步回答“有限的当前认知表面怎样连接一个远比自身庞大的多尺度智能系统”。后续关于自我感知、认知控制和运行态治理的机制，都将建立在这一认知总线之上。
\end{chapteroverviewbox}

\section{从工作记忆到认知总线：理论对象的重新定义}

如果我们仍然沿用传统计算机体系中的直觉，工作记忆很容易被误解为一种容量较小、访问速度较高的内部存储。这样的比喻虽然能够帮助我们理解 ``short-term storage'' 的部分现象，却不足以描述一个持续智能体真实的运行结构。原因在于，一个 Agent 的当前认知并不仅仅是在保存若干信息，它还必须同时完成目标维持、假设竞争、情境解释、记忆召回、身体状态接入、行动规划、工具调用以及自我状态进入等功能。换句话说，$h(t)$ 中真正稀缺的并不是字节数量，而是\textbf{能够被同时保持为全局可访问关系结构的认知承载能力}。

因此，我们首先把工作记忆重新定义为一个随时间变化的结构状态：
\[
h(t)=
\left(
V_h(t),
E_h(t),
W_h(t),
A_h(t)
\right),
\]
其中 $V_h(t)$ 表示当前显式激活的 Agent-CSO 集合，$E_h(t)$ 表示这些结构之间当前显式维持的关系集合，$W_h(t)$ 表示关系强度、优先级或有效注意权重，而 $A_h(t)$ 表示当前具有行动相关性的激活状态。这个定义意味着 $h(t)$ 不是一个线性 token 序列，也不是一个静态 KV-cache，而是一张规模有限、持续变化的当前认知关系图。

设整个 Agent 生命周期中潜在可访问的认知结构集合为
\[
\mathcal C_A(t),
\]
显然一般有
\[
|V_h(t)|\ll |\mathcal C_A(t)|.
\]
更重要的是，即使 $|V_h|$ 较小，关系数量也可能快速增长。若所有节点之间都形成潜在关系，则最坏情况下：
\[
|E_h|
=
O(|V_h|^2).
\]
因此真正限制工作记忆的并不仅是当前项目数，而是活跃 Agent-CSO 之间需要被同步维持的有效关系复杂度。这正是上一章有限容量定理的结构来源。为了使智能体维持可运行状态，我们需要：
\[
C_h(t)
=
\Phi
\left(
V_h(t),
E_h(t),
W_h(t)
\right)
\leq C_h^{\max},
\]
其中 $C_h(t)$ 是当前工作记忆中的有效结构负载，而 $C_h^{\max}$ 是由底层模型、表征压缩率、计算资源、当前认知温度和其他实现条件共同决定的可维持上界。

从这里我们就可以看到，把 $h(t)$ 称为“工作记忆”只描述了它的存储侧功能，而把它称为“认知总线”则强调了它真正重要的另一面：\textbf{不同时间尺度、不同载体和不同功能区域中的认知结构，只有在获得进入 $h(t)$ 的显式激活资格之后，才进入当前全局推理、目标竞争和行动闭合过程。}

因此，$h(t)$ 的更准确描述应当是：
\[
\boxed{
h(t)
=
\text{Sparse Explicit Cross-Scale Cognitive Surface}
}
\]
即“稀疏的跨尺度显式认知表面”。

这里的 ``Sparse'' 表示只有极少部分生命周期结构被当前激活；``Explicit'' 表示这些结构当前能够参与全局组合、比较和控制；``Cross-Scale'' 表示它连接比自己更快和更慢的动力层；``Surface'' 则表示它不是整个智能本体，而只是整个智能系统在当前时刻被暴露出来的一个截面。

这一重新定义具有非常重要的认识论后果。一个成熟 Agent 的绝大多数能力实际上不应该持续存在于 $h(t)$ 中。已经成熟的动作技能可以存在于更低层快速闭环；稳定世界知识可以存在于 $\Theta$；完整经历可以存在于 $E$；长期自我结构主要由 $\Psi$ 承载；程序、文件、数据库和工具则可能位于外部载体。于是：
\[
\boxed{
MostAgentDynamics\notin h(t)
}
\]
并不是工作记忆的缺陷，反而是成熟智能成立的必要条件。

我们甚至可以把智能成熟程度与显式认知依赖建立一种概念关系。设系统成功完成某类行为所需的平均显式工作记忆介入量为 $\mathcal I_h$，系统完成该行为的有效能力为 $\mathcal K$，则可以定义一个概念性的显式依赖比：
\[
R_h
=
\frac{\mathcal I_h}{\mathcal K}.
\]
当一个能力从陌生任务逐步变成熟练技能时，通常希望：
\[
R_h\downarrow.
\]
这说明学习的成熟并不是让工作记忆越来越庞大，而是让反复发生的认知结构逐步迁移到更稳定、更局部、更低成本的承载位置，使 $h(t)$ 重新释放给真正新颖、冲突和需要跨尺度协调的问题。

由此，本章的第一个核心原则可以写为：
\begin{proposition*}
工作记忆的价值不在于容纳最多内容，而在于成为有限认知容量下最关键结构的当前公共工作面。
\end{proposition*}

\section{六类基本通路：认知总线连接什么}

既然 $h(t)$ 是认知总线，我们接下来就必须回答一个比“工作记忆里放什么”更严格的问题：\textbf{哪些不同性质的结构必须通过这条总线发生交汇？} 从此前整个 AgentOS 原理论可以归纳出六类最基本的信息与认知结构通路：
\[
World\rightarrow h,
\]
\[
E\rightarrow h,
\]
\[
\Theta\rightarrow h,
\]
\[
Self\rightarrow h,
\]
\[
h\rightarrow Action,
\]
\[
h\rightarrow E/Me.
\]

需要强调的是，这六条箭头并不表示六条固定的物理数据线，而表示六种不同来源和不同时间方向的结构进入或离开当前显式认知表面的动力关系。

第一类是
\[
World\rightarrow h.
\]
现实世界在任何时刻都包含远超工作记忆容量的状态变化：
\[
X_W(t)\in\mathcal X_W.
\]
Agent 不可能把 $X_W(t)$ 完整映射到 $h(t)$。真实过程更接近：
\[
X_W(t)
\xrightarrow{\mathcal O}
Z_W(t)
\xrightarrow{\mathcal A}
h_W(t),
\]
其中 $\mathcal O$ 是观测与初步表示过程，$\mathcal A$ 是准入与显式激活过程。最终进入 $h$ 的只是：
\[
h_W(t)\subset Z_W(t),
\]
即当前与目标、风险、新颖性、身体状态或认知冲突相关的有限结构。因此，现实进入工作记忆不是复制，而是选择性结构投影。

第二类是：
\[
E\rightarrow h.
\]
$E$ 是情景轨迹记忆，它保存的不只是事实，而是经历在时间中的结构。当前任务可能触发某段过去：
\[
E_k
=
\{e(\tau)\mid \tau\in[t_0,t_1]\}.
\]
通过情景匹配：
\[
R_E:
(h(t),E)\rightarrow \hat E_k,
\]
系统得到与当前情境相关的经验片段，但这些片段仍不能全部进入当前显式面，因此还需要：
\[
\Pi_h(\hat E_k)
\rightarrow h_E(t).
\]
这使得回忆成为一种“过去结构重新进入现在”的过程。我们由此也能看出，$h$ 不只是现在，它还是过去能够重新影响现在的入口。

第三类是：
\[
\Theta\rightarrow h.
\]
$\Theta$ 表示较慢的结构记忆和生成结构。它所提供的不是某个具体经历，而是：
\[
Rule,\quad Schema,\quad Skill,\quad WorldModel,\quad Prior,\quad PredictionStructure.
\]
可以概念性表示为：
\[
h_\Theta(t)
=
G_\Theta
\left(
h(t),
\Theta
\right),
\]
其中 $G_\Theta$ 根据当前认知状态激活相关结构。这样，当前工作记忆中的某个对象并非孤立存在，而是会自动获得由 $\Theta$ 提供的解释框架、可能关系和预测约束。

第四类是：
\[
Self\rightarrow h.
\]
这里的 Self 主要指由 $\Psi$ 所形成的长期自我结构对当前认知的显式投影。严格说，$\Psi$ 并不会全部进入 $h$；它主要以慢约束场存在。但在某些时刻，自身状态、长期偏好、身份相关冲突或者价值判断需要显式进入当前工作面，因此形成：
\[
h_{\mathrm{self}}(t)
=
\Pi_{\mathrm{self}}
\left(
\Psi,E,h
\right).
\]
这条通路非常重要，因为没有它，一个系统即使能够调用世界知识和过去经验，也仍然只是一个无主体归属的计算过程。只有当“这个事实”“这段经历”“这个目标”进一步被表达成“它对当前这个我意味着什么”，认知才真正获得主体坐标。

第五类是：
\[
h\rightarrow Action.
\]
工作记忆不是内向封闭空间，而必须通过行动使认知进入现实。设当前显式状态为 $h(t)$，行动候选集合为：
\[
\mathcal A_t
=
\{a_1,\dots,a_m\}.
\]
则行动生成可写为：
\[
a_t^*
=
\arg\max_{a\in\mathcal A_t}
J(a\mid h(t),\Theta,\Psi),
\]
其中 $J$ 不必是单一标量奖励，而可以代表多目标约束下的行动选择函数。关键在于：只有进入当前显式认知的结构才能直接参与这一层级的全局行动决定。那些未进入 $h$ 的低层能力仍然可以执行动作，但它们通常作为局部闭环完成，而非形成全局显式决策。

第六类是：
\[
h\rightarrow E/Me.
\]
当前认知并不只是消费旧结构，它还不断形成未来结构。当前显式处理完成以后，一部分内容进入：
\[
h\rightarrow E,
\]
成为新的情景经验；另一部分可能被外化到：
\[
Me
=
\{
Skill,File,DB,Program,Tool,Environment,\ldots
\},
\]
形成：
\[
h\rightarrow Me.
\]
因此，$h$ 同时是记忆和外部结构的写入口。

六类通路组合以后，$h(t)$ 的中心位置可以写成：
\[
\boxed{
h(t)
=
\mathcal B_C
\left(
World,
E,
\Theta,
\Psi,
Me,
Action
\right)
}
\]
其中 $\mathcal B_C$ 表示 Cognitive Bus，而不是简单数据求和。

由此我们得到第二个核心原则：
\begin{proposition*}
工作记忆不是一种来源的信息空间，而是不同来源、不同时间尺度和不同承载位置的认知结构在当前时刻发生显式耦合的公共接口。
\end{proposition*}

\section{进入工作记忆意味着什么：Agent-CSO 的显式激活}

为了避免把“进入工作记忆”理解为数据复制，我们必须回到前面已经建立的 CSO 与 Agent-CSO 理论。设现实或认知本体中的某个结构为：
\[
C\in\mathcal C,
\]
Agent 对其具体承载为：
\[
AgentCSO_A(C,t)
=
(C,R,F,\tau,B,S).
\]
其中 $R$ 表示当前表示形式，$F$ 表示功能承载位置，$\tau$ 表示有效时间尺度，$B$ 表示承载边界，而 $S$ 表示激活状态。

那么所谓某个认知结构“进入 $h$”，更精确地应该写成：
\[
S:
Latent
\rightarrow
ExplicitActive,
\]
以及：
\[
F:
F_i
\rightarrow
F_h
\]
或者更一般地：
\[
\mu_t(C)
\ni
F_h.
\]
这意味着同一个 CSO 不一定离开原有承载位置；更常见的情况是，它在保留原有长期承载的同时获得了一个当前工作记忆投影。例如，某个长期知识结构仍然存在于 $\Theta$，但其中少量与当前任务相关的关系被激活进入 $h$：
\[
\Theta_C
\rightarrow
h_C.
\]
同样，一段完整经历仍存在于 $E$，工作记忆只得到它的当前相关压缩：
\[
E_k
\rightarrow
\operatorname{Summary}(E_k)
\rightarrow
h_E.
\]

因此，工作记忆是一个\textbf{激活空间}，而不是结构唯一的“住所”。

这一点对于整个 AgentOS 理论极为重要。因为如果我们把进入 $h$ 理解为完整搬运，就会自然产生一个错误架构：所有记忆、所有工具状态、所有身体状态都必须被展开到一个巨大上下文中。随着系统持续运行，这种方法必然导致：
\[
ContextSize(t)\uparrow,
\]
进而使：
\[
SearchCost\uparrow,
\qquad
Interference\uparrow,
\qquad
Latency\uparrow,
\]
最终破坏工作记忆的有效结构容量。

更合理的认知机制应当是：
\[
\boxed{
LatentStructure
\rightarrow
SelectiveActivation
\rightarrow
SparseExplicitProjection
}
\]
而不是：
\[
AllStructure\rightarrow Context.
\]

我们可以进一步定义某个 Agent-CSO $C_i$ 在时刻 $t$ 获得工作记忆激活资格的概率：
\[
P_i^{h}(t)
=
\sigma
\left[
\alpha_rR_i
+
\alpha_gG_i
+
\alpha_nN_i
+
\alpha_eE_i
+
\alpha_sS_i
-
\alpha_cC_i
\right],
\]
其中：

$R_i$ 表示当前相关性；

$G_i$ 表示目标相关性；

$N_i$ 表示新颖性；

$E_i$ 表示异常或残差强度；

$S_i$ 表示自我/价值相关性；

$C_i$ 表示进入当前显式空间的结构成本；

$\sigma$ 为归一化准入函数。

这里我们并不要求这一公式直接成为工程实现，而是用它说明一个重要原理：\textbf{显式认知是一种竞争性资源分配，而不是默认状态。}

于是可以定义当前工作记忆：
\[
V_h(t)
=
\left\{
C_i
\mid
P_i^{h}(t)>\theta_h
\right\}.
\]
更现实地，因为有限容量约束，还需要满足：
\[
\sum_{C_i\in V_h(t)}
c_i(t)
+
\sum_{(i,j)\in E_h(t)}
c_{ij}(t)
\leq
C_h^{\max}.
\]
这意味着即使多个结构都具有进入理由，最终仍可能因为全局容量限制发生竞争。

这也使“注意”获得一个比普通 token 权重更深的定义：
\[
\boxed{
Attention
=
PolicyForPromotingAgentCSOIntoExplicitCrossScaleProcessing
}
\]
即注意力本质上是某个结构获得更高层显式处理权的策略。这里的注意不同于 Transformer 内部 attention operator；前者属于认知架构层，后者属于模型实现层，两者不能混为一谈。

进一步，当某个 Agent-CSO 从 $h$ 中退出，也不意味着被删除。其状态可能发生：
\[
ExplicitActive
\rightarrow
LatentAvailable,
\]
或者：
\[
h\rightarrow E,
\]
或者：
\[
h\rightarrow\Theta,
\]
或者：
\[
h\rightarrow Me.
\]
因此工作记忆真正管理的是：
\begin{statementbox}
Explicitness
\end{statementbox}
而不是：
\[
Existence.
\]

从这个角度出发，我们可以得到第三个重要原则：
\begin{proposition*}
进入工作记忆不是“存在”，而是获得当前全局显式认知权；退出工作记忆也不是“消失”，而是认知结构重新回到其他时间尺度或承载位置。
\end{proposition*}

\section{h(t) 与 E、$\Theta$、$\Psi$ 的多时间尺度耦合}

如果把工作记忆视为跨尺度显式表面，那么它最核心的理论任务之一，就是在不同时间尺度之间建立可控耦合。我们已经有：
\[
\tau_h
\ll
\tau_E
\ll
\tau_\Theta
\ll
\tau_\Psi.
\]
这意味着 $h$ 变化最快，$E$ 较慢，$\Theta$ 更慢，而 $\Psi$ 是更长期的身份与价值约束结构。这样的时间尺度分离不只是记忆分类，而意味着不同层级在动力学上承担不同责任。

可以建立一个最小耦合模型：
\[
\dot h
=
F_h
\left(
h,
u_W,
r_E,
r_\Theta,
r_\Psi
\right),
\]
其中：

$u_W$ 是当前世界输入；

$r_E$ 是情景记忆召回；

$r_\Theta$ 是结构记忆激活；

$r_\Psi$ 是自我约束投影。

同时：
\[
\dot E
=
\epsilon_E
F_E(E,h,\Theta),
\]
\[
\dot\Theta
=
\epsilon_\Theta
F_\Theta(\Theta,E,h),
\]
\[
\dot\Psi
=
\epsilon_\Psi
F_\Psi(\Psi,E,\Theta),
\]
满足：
\[
1
\gg
\epsilon_E
\gg
\epsilon_\Theta
\gg
\epsilon_\Psi.
\]

这套方程表达了一个非常关键的设计原则：\textbf{快变量可以频繁变化，慢变量只能吸收经过时间检验的结构残差。}

如果每一次 $h$ 的短暂变化都立即修改 $\Theta$ 或 $\Psi$，系统会产生：
\[
FastNoise
\rightarrow
SlowStructure,
\]
最终导致长期结构被瞬时噪声污染。反之，如果慢变量完全不受 $h$ 影响，则系统无法持续学习。因此更合理的方式是通过时间积分或重复证据形成慢变量更新：
\[
\Delta\Theta
\propto
\int_{t_0}^{t_1}
w_E(\tau)\,
\phi(h(\tau))
\,d\tau,
\]
而：
\[
\Delta\Psi
\propto
\int_{T_0}^{T_1}
w_\Psi(\tau)
\,
\chi(E(\tau),\Theta(\tau))
\,d\tau,
\]
其中通常有：
\[
T_1-T_0
\gg
t_1-t_0.
\]

这意味着 $h$ 对慢层的影响不是直接写入，而是：
\[
h
\rightarrow
RepeatedExperience
\rightarrow
E
\rightarrow
StableStructure
\rightarrow
\Theta
\rightarrow
LongTermSelfInterpretation
\rightarrow
\Psi.
\]

反向影响则不同。慢结构并不需要完全展开进入 $h$，它们可以通过约束方式影响当前认知。例如：
\[
\Theta
\rightarrow
Prior(h),
\]
\[
\Psi
\rightarrow
Preference(h).
\]
因此，一个成熟系统中大量慢结构对快层的影响应该是“背景性的”，而不是持续显式占据工作记忆。这可以概括成：
\[
\boxed{
SlowConstraint\neq ExplicitOccupation
}
\]

这一点特别重要。比如一个成熟 Agent 已经形成“保护人类安全”这样的长期结构，它不应该在每一次动作中都把完整价值论证展开到 $h$。更合理的是，这种慢结构定义当前可行空间：
\[
\mathcal A_{\mathrm{feasible}}
=
\left\{
a
\mid
g_\Psi(a)\leq 0
\right\},
\]
而 $h$ 只在这个约束区域中搜索具体行动。只有当发生价值冲突、边界事件或异常残差时，相关慢结构才需要被显式提升：
\[
\Psi
\rightarrow
h.
\]

因此，不同时间尺度之间的关系并不是：
\[
SlowLayer\to FastLayer
\]
持续输出详细命令，而是：
\[
\boxed{
Higher/SlowerLayer
\rightarrow
FeasibleRegion,\ Prior,\ Gain,\ Boundary
}
\]
而快层在该区域中保留局部自治。

这与我们更一般的多尺度智能原则一致：
\[
\boxed{
HigherLevelSetsFeasibleRegion,\qquad
LowerLevelOptimizesInsideIt
}
\]

工作记忆 $h$ 因此具有一个非常特殊的位置：它既不是最快层，也不是最慢层，而是\textbf{多尺度结构需要发生显式协商时形成的中介面}。在正常情况下，低层熟练能力应保持局部；长期结构应保持背景性；只有新颖、冲突、异常和跨层协调问题才需要占据 $h$。

由此得到本节核心：
\begin{statementbox}
h(t) 的功能不是把所有时间尺度统一到一个速度，而是在必要时为不同时间尺度提供有限的显式耦合窗口。
\end{statementbox}

\section{h(t) 与世界：现实不是连续变成思想}

讨论工作记忆时，一个常见误解是认为传感器不断获得的数据应该不断流入认知核心。对于真正的具身智能体，这种设计既不现实，也违背工作记忆有限性。世界具有远高于当前显式认知的状态维度：
\[
\dim(X_W)
\gg
\dim(h).
\]
因此，世界与工作记忆之间必然存在大量局部闭环、过滤、压缩和事件化过程。

可以将外部现实到工作记忆的映射写成：
\[
X_W(t)
\xrightarrow{\mathcal P}
Z(t)
\xrightarrow{\mathcal F}
e(t)
\xrightarrow{\mathcal G}
h_W(t),
\]
其中：

$\mathcal P$ 表示感知与低层表示；

$\mathcal F$ 表示局部结构提取、预测和残差检测；

$e(t)$ 表示具有认知意义的事件；

$\mathcal G$ 表示将事件提升到显式认知空间。

因此：
\begin{statementbox}
RealityDoesNotContinuouslyBecomeThought
\end{statementbox}
更准确的是：
\begin{statementbox}
RelevantDeviationBecomesThought.
\end{statementbox}

设低层预测：
\[
\hat x(t+\Delta)
=
P(x(t)),
\]
实际观测为：
\[
x(t+\Delta),
\]
则残差：
\[
\varepsilon(t)
=
x(t+\Delta)-\hat x(t+\Delta).
\]
如果：
\[
\|\varepsilon(t)\|<\theta_\varepsilon,
\]
并且风险、新颖度与目标相关度都较低，则局部层可以自行闭合，不必提升至 $h$。只有当：
\[
\Gamma(t)
=
F
\left(
\|\varepsilon\|,
Risk,
Novelty,
Persistence,
GoalRelevance
\right)
\geq
\theta_h
\]
时，相关结构才进入：
\[
World\rightarrow h.
\]

这使得 h 的“世界输入”本质上是一种\textbf{事件触发式跨尺度升级}，而不是连续带宽灌入。它解释了为什么一个成熟的人在走路时不会持续显式计算每一步的肌肉控制，也解释了为什么未来 AgentOS 不应让大模型逐帧处理所有传感器数据。

从认知角度说，世界中大多数正常变化应当由：
\[
LocalPrediction
+
LocalControl
\]
处理；工作记忆主要处理：
\[
Novelty
+
Conflict
+
Uncertainty
+
GoalTransition
+
Residual.
\]

这也使我们能够重新定义“注意力”。注意并不意味着系统此刻唯一处理某个对象，而意味着某个对象的局部动态已经获得足够重要性，需要提升到显式全局表面：
\[
\boxed{
Attention
=
CrossScalePromotion.
}
\]

世界与 h 之间还存在一个更深的认识论问题：进入 h 的永远不是真实世界本身，而是 Agent 根据其观测能力、已有结构与当前目标形成的 Agent-CSO。因此：
\[
h_W(t)
\neq
X_W(t).
\]
严格地说：
\[
h_W(t)
=
\Pi_A
\left[
O_A(X_W(t))
\right].
\]
这意味着工作记忆天然具有认识论不完备性。

于是任何当前显式信念都应允许携带：
\[
Confidence,
Source,
Timestamp,
Observed/Derived/PredictedStatus.
\]
这为后续工程中的可验证语义状态提供理论基础，但本章只需要确立一个原则：
\[
\boxed{
Representation\neq Reality.
}
\]

因此，当行动从 h 发出以后，系统不能把预测的成功结果直接写回认知总线，而必须等待现实重新进入：
\[
h
\rightarrow
Action
\rightarrow
Reality
\rightarrow
Observation
\rightarrow
h'.
\]
这里：
\[
ActionIssued
\neq
ActionSucceeded.
\]
这一闭环保证了认知总线不会逐渐退化成内部自洽但与现实脱离的生成空间。

由此本节可凝练为：
\begin{statementbox}
世界不是持续进入工作记忆；世界只有在局部闭环无法继续可靠处理、或其变化具有足够认知价值时，才以结构化事件的形式进入工作记忆。
\end{statementbox}

\section{h(t) 与行动：显式认知必须重新闭合到现实}

如果只讨论：
\[
World\rightarrow h,
\]
那么工作记忆仍然只是一个观察系统。真正的智能必须形成：
\[
World
\rightarrow
h
\rightarrow
Action
\rightarrow
World'.
\]
因此，$h(t)$ 的第二个本质不是记忆，而是\textbf{行动选择的显式协调面}。

设当前 h 中包含：
\[
h(t)=
\{
G_t,
B_t,
C_t,
P_t,
U_t,
S_t
\},
\]
分别表示当前目标、信念、约束、计划、不确定性以及自我状态。行动选择不是由其中某一个变量独立决定，而是这些当前显式结构形成的联合约束问题。

可以写成：
\[
a_t^*
=
\arg\min_{a\in\mathcal A}
\Big[
L_G(a)
+
\lambda_RL_R(a)
+
\lambda_UL_U(a)
+
\lambda_SL_S(a)
\Big],
\]
其中：

$L_G$ 表示目标偏离；

$L_R$ 表示风险；

$L_U$ 表示不确定性；

$L_S$ 表示与慢层自我和长期结构的不一致。

然而，真正的 AgentOS 不应要求 h 直接生成所有低级执行细节。因此更合理的是：
\[
h
\rightarrow
Intent/Plan/Reference
\rightarrow
LowerLoop.
\]
也就是说，h 的行动输出应主要表示：
\[
WhatToDo
\]
以及：
\[
WhatConstraintsMustHold,
\]
而不是无限向下展开：
\[
HowEveryPrimitiveMotionOccurs.
\]

从纯认知理论角度，我们可以把行动看成工作记忆把内部结构施加到现实上的投影：
\[
\mathcal A:
h(t)
\rightarrow
u(t).
\]
现实动力学：
\[
\dot x_W
=
F_W(x_W,u).
\]
随后新的现实状态：
\[
x_W(t+\Delta)
\]
重新经过观测进入 Agent。于是形成：
\[
\boxed{
h
\rightarrow
Action
\rightarrow
World
\rightarrow
h'
}
\]
这一闭合。

这个闭环具有两个重要含义。

第一，工作记忆中的推理必须具有\textbf{现实可闭合性}。如果某个长期推理不断占据 h 却既不能产生行动，也不能形成可验证内部结构，那么它将逐渐变成一种认知驻留。对于容量有限的 h，这是非常昂贵的。因此成熟系统必须努力使当前认知过程趋向：
\[
Goal
\rightarrow
Inference
\rightarrow
Decision
\rightarrow
Action/Commit
\rightarrow
Closure.
\]

第二，行动本身也是认知的一部分。因为某些世界结构只有通过行动才能被识别。例如：
\[
ObservationAlone
\]
可能不能区分两个假设，而一个主动试探：
\[
a_{\mathrm{probe}}
\]
可以改变世界状态，使两个假设产生不同后果。因此：
\[
Action
\]
同时具有控制和信息获取功能。

可将主动信息行为定义为：
\[
a^*
=
\arg\max_a
\left[
\mathbb E(U(a))
+
\beta
\mathcal I
\left(
Observation_{t+1};
Hypothesis
\mid a
\right)
\right].
\]
这里的第二项表示行动带来的信息收益。

所以：
\begin{statementbox}
Thinking
\end{statementbox}
与：
\begin{statementbox}
Acting
\end{statementbox}
并不是两个彼此完全分离的过程。工作记忆中的某些不确定性只有通过现实交互才能消解。于是：
\[
h\rightarrow Action
\]
实际上也是一种：
\[
h\rightarrow BetterFutureh.
\]

这使工作记忆获得了更加完整的动力学定义：
\[
\boxed{
h(t)
=
\text{当前 Agent-CSO 在预测、选择、行动和现实验证之间发生显式闭合的有限结构面。}
}
\]

\section{h(t) 与外部认知载体：显式空间不是全部内部空间}

工作记忆有限性不仅要求系统把结构迁移到不同内部记忆层，也必然要求智能体利用外部世界承担部分认知结构。在早期讨论中，我们把这一过程称为“认知复杂度卸载”；在后来 CSO 本体框架中，更严格的理解是：
\begin{statementbox}
AgentCSOCarrierMigration.
\end{statementbox}

设某个当前 Agent-CSO：
\[
C_i@h.
\]
如果它持续占据工作记忆，但其内容具有稳定、可恢复、重复使用的性质，那么继续保持它处于显式工作状态是不经济的。系统可以执行：
\[
C_i@h
\rightarrow
C_i@E,
\]
或者：
\[
C_i@h
\rightarrow
C_i@\Theta,
\]
也可以：
\[
C_i@h
\rightarrow
C_i@Me.
\]

这里：
\[
Me
=
\{
File,
Database,
Script,
Program,
Skill,
Tool,
Environment,
OtherAgent,
\ldots
\}.
\]
$Me$ 不表示所有外部资源，而表示已经被 Agent 建立稳定认知索引、并可以在未来重新调用的外部承载结构。

因此，一个文件只有在满足：
\[
Index(File)
+
Retrievability(File)
+
SemanticRelation(File,Agent)
\]
时，才真正成为 Agent 外部记忆的一部分。

例如，当前工作记忆中维持一个复杂计算过程：
\[
h:
\{x_1,x_2,\dots,x_n\}
\]
如果这些中间结构持续消耗 h 容量，可以将它们写入文件或程序：
\[
h
\rightarrow
File/Program,
\]
然后在 h 中只保留：
\[
Pointer+
Meaning+
Status.
\]
于是工作记忆的负载从：
\[
C_{\mathrm{full}}
\]
下降为：
\[
C_{\mathrm{index}}
\ll
C_{\mathrm{full}}.
\]

这正是认知卸载能够增加实际问题规模的关键原因。它并没有让工作记忆容量变大，而是减少了每个任务结构需要在当前显式空间中保持的尺寸。

从信息结构角度，可以定义压缩比：
\[
\kappa_i
=
\frac{C_{explicit}(C_i)}
{C_{index}(C_i)}.
\]
若：
\[
\kappa_i\gg1,
\]
则外化具有很高价值。

但外化也存在恢复代价。设外部化代价：
\[
C_{off},
\]
重新加载代价：
\[
C_{recall},
\]
外部失败概率：
\[
p_f,
\]
则一个结构是否值得卸载，需要比较：
\[
Benefit_{off}
=
Saving_h
-
C_{off}
-
\mathbb E[C_{recall}]
-
\lambda p_f.
\]
所以“把东西放外部”并非无条件有利。

更重要的是，外部结构可能被环境本身承载。例如，一个 Agent 每次进入房间都需要重新计算某个物品放置规则，那么它可以直接修改环境，使规则成为环境结构：
\[
InternalRule
\rightarrow
EnvironmentalConstraint.
\]
此后现实本身帮助 Agent 降低认知负担。这正是我们早期提出：
\begin{statementbox}
EnvironmentAsExternalMemory
\end{statementbox}
的根本含义。

因此：
\[
h\rightarrow Me
\]
不是一个附加功能，而是有限认知系统扩大问题求解空间的重要机制。

反过来：
\[
Me\rightarrow h
\]
则表示外部结构被重新激活。例如读取文件、调用 Skill、向其他 Agent 查询结果，本质上都是：
\[
ExternalAgentCSO
\rightarrow
CurrentExplicitAgentCSO.
\]

由此，我们可以把工作记忆理解成内部世界与外部认知世界之间的双向交换口：
\[
\boxed{
InternalStructure
\leftrightarrow
h
\leftrightarrow
ExternalStructure.
}
\]

这也使“认知边界”不再等同于身体边界。一个成熟 Agent 的有效认知范围可以跨越内部模型、文件系统、工具、环境与协作主体。真正持续的智能并不是把所有知识都装进自己，而是能够稳定知道：
\[
WhatMustRemainInternal,
\]
\[
WhatCanLiveOutside,
\]
以及：
\[
HowToRecoverItWhenNeeded.
\]

\section{工作记忆的拓扑：局部关系图而不是线性上下文}

为了进一步提高理论精度，我们需要讨论 $h(t)$ 的内部几何结构。当前许多语言模型系统事实上把工作区表示成序列：
\[
h(t)
=
(x_1,x_2,\dots,x_n).
\]
然而，从 AgentOS 的角度看，这种线性形式只是具体实现之一，工作记忆的本体更接近动态关系图。

定义：
\[
\mathcal G_h(t)
=
\left(
V_h(t),
E_h(t)
\right).
\]
其中节点可以包括：

当前目标；

子目标；

实体；

假设；

关键状态；

约束；

计划步骤；

当前世界对象；

自我状态；

工具状态；

开放问题；

关键中间变量。

边则可以表示：

因果；

时序；

依赖；

冲突；

支持；

归属；

目标关系；

空间关系；

资源竞争；

条件关系。

于是一个当前任务真正的认知复杂度可能更接近：
\[
C_h
=
\alpha |V_h|
+
\beta |E_h|
+
\gamma H(E_h)
+
\delta D_h,
\]
其中 $H(E_h)$ 可以表示关系不确定性，$D_h$ 表示跨节点耦合深度。

这样我们就能解释一个现实中非常常见的现象：两个任务即使包含相同数量的信息，其工作记忆负担也可能完全不同。若第一个任务由十个几乎独立的对象组成：
\[
|V|=10,\qquad |E|\approx 5,
\]
而第二个任务同样有十个对象，却存在大量交叉约束：
\[
|V|=10,\qquad |E|\approx40,
\]
那么后者的认知负担显然更高。

因此：
\[
\boxed{
WorkingMemoryLoad
\neq
InformationAmount.
}
\]
更加准确的是：
\[
\boxed{
WorkingMemoryLoad
\approx
ActiveRelationalStructure.
}
\]

这也是为什么简单扩展 context length 不能从根本上解决复杂推理问题。上下文长度主要扩大：
\[
StorageCapacity,
\]
而工作记忆真正要求的是：
\[
CoherentActiveRelationCapacity.
\]
如果增加大量未被结构化的信息：
\[
n\uparrow
\]
反而可能使相关结构提取难度和干扰增加。

因此，AgentOS 必须允许 h 形成多级折叠表示。例如一个复杂子问题：
\[
\mathcal G_{sub}
\]
经过结构压缩后形成一个高阶节点：
\[
v_{macro}
=
Fold(\mathcal G_{sub}).
\]
于是：
\[
\mathcal G_h
\rightarrow
\mathcal G_h',
\]
并满足：
\[
|V_h'|<|V_h|,
\qquad
|E_h'|<|E_h|,
\]
同时尽量保持：
\[
TaskRelevantInformation
\]
不显著下降。

这就是所谓认知折叠：
\[
\boxed{
Fold:
DetailedStructure
\rightarrow
ActionRelevantMacroStructure.
}
\]

它和传统摘要有所区别。摘要通常压缩文本，而 Fold 压缩的是\textbf{当前决策所需的关系结构}。

进一步，工作记忆还必须允许局部子图暂时退出：
\[
\mathcal G_{branch}
\subset
\mathcal G_h
\]
若当前不需要继续推理，则：
\[
\mathcal G_{branch}
\rightarrow
PausedState,
\]
并只留下恢复句柄：
\[
Handle(\mathcal G_{branch}).
\]
未来需要时：
\[
Handle
\rightarrow
Restore(\mathcal G_{branch}).
\]

这种设计与操作系统中的进程挂起存在形式类比，但需要注意：认知线程不是 CPU 线程。一个认知线程是一组相互依赖的 Agent-CSO 及其未闭合状态。它的暂停和恢复要求保持语义连续性，而不仅是寄存器现场。

从图论角度，可以定义当前工作记忆中某个节点的中心性：
\[
Cent(v_i)
=
f(
Degree,
Betweenness,
GoalRelevance,
Residual
).
\]
高中心性节点更可能成为当前认知枢纽。例如当前目标、关键约束、核心冲突等往往形成：
\[
Hub_h.
\]
但如果某个节点长期拥有过高中心性，则可能导致所有认知都围绕单一点展开，形成认知僵化；反过来，如果图结构过度均匀且缺乏稳定中心，则可能出现目标漂移。因此 h 的良好状态并不是最大连接，而是\textbf{有限容量下的适度集中与可重组性}。

这为后续工作记忆运行相态提供了自然基础，但即使不引入任何相态概念，我们已经可以得到一个非常明确的工程结论：
\begin{conclusion}
AgentOS 的工作记忆应首先被实现为具有可折叠、可挂起、可恢复、可重构关系结构的动态认知图，而不是一个无限增长的文本上下文。
\end{conclusion}

\section{认知总线上的闭环完整性：目标、假设、行动与验证}

工作记忆成为认知总线以后，还必须解决一个容易被忽视的问题：一个认知结构不仅需要进入 h，还必须具有生命周期。否则工作记忆会不断积累未完成的目标、未验证的假设和悬挂的中间状态。

因此我们定义一个显式认知对象：
\[
q_i(t)
=
\left(
G_i,
H_i,
P_i,
A_i,
V_i,
S_i
\right),
\]
其中：

$G_i$ 是目标；

$H_i$ 是当前假设；

$P_i$ 是过程/计划状态；

$A_i$ 是行动或待执行操作；

$V_i$ 是验证条件；

$S_i$ 是认知线程状态。

一个健康的认知线程应该倾向于形成：
\[
Goal
\rightarrow
Hypothesis
\rightarrow
Plan
\rightarrow
Action
\rightarrow
Verification
\rightarrow
Closure.
\]
我们可以把它记作：
\[
\mathcal L_q.
\]

如果行动结果满足：
\[
V_i(World_{t+\Delta})=1,
\]
则线程可以进入：
\[
Closed.
\]
如果：
\[
V_i=0,
\]
则产生残差：
\[
\varepsilon_i,
\]
重新进入：
\[
HypothesisRevision
\]
或：
\[
PlanRevision.
\]

这种闭环完整性对有限工作记忆极为关键。因为一个长期未闭合线程具有持续占用：
\[
C_i(t)>0.
\]
如果系统累积：
\[
N_{open}\uparrow,
\]
则：
\[
C_h
=
\sum_{i=1}^{N_{open}}C_i
+
C_{\mathrm{cross}}
\]
可能迅速接近容量极限。

因此需要定义一个开放认知债务：
\[
D_{\mathrm{open}}(t)
=
\sum_i
w_i
\mathbb I[S_i\neq Closed].
\]
如果：
\[
D_{\mathrm{open}}\gg D_{\max},
\]
即使每个问题都不复杂，系统也可能因为大量悬挂线程进入失稳。

这解释了为什么一个智能体可能不是“不会推理”，而是“同时承载了太多未完成的认知责任”。从 AgentOS 的角度看，任务完成能力的一部分就是：
\begin{statementbox}
CognitiveClosureManagement.
\end{statementbox}

进一步，一个线程并不一定需要在 h 中等待现实结果。例如外部任务需要数小时才能完成，如果持续保留在工作记忆：
\[
Task_i@h
\]
显然非常浪费。更合理的是：
\[
Task_i
\rightarrow
ExternalPendingState,
\]
只保留：
\[
Trigger_i.
\]
当结果出现：
\[
Event_i
\]
再重新提升：
\[
Event_i\rightarrow h.
\]

这意味着认知总线应天然是事件驱动，而非阻塞式的。一个长期 Agent 不应因为某个任务等待而使整个主体停止。这与我们此前形成的持续主体原则是一致的：任务可以等待，主体不能被任务困住。

因此，认知总线上的过程更接近：
\[
Event
\rightarrow
Activate
\rightarrow
Process
\rightarrow
Commit/Delegate
\rightarrow
Suspend/Close
\]
而不是传统函数调用：
\[
Call
\rightarrow
Wait
\rightarrow
Return.
\]

这一点将对未来 AgentOS 的进程模型产生直接工程意义，但在本章，我们只需要确立其认知原理：\textbf{显式认知资源必须通过闭环、暂停或外化不断被释放，否则有限工作记忆最终一定被历史未闭合状态淹没。}

由此可以写出：
\begin{statementbox}
工作记忆不仅管理“什么正在被思考”，还必须管理“什么已经完成、什么需要等待、什么可以退出以及什么必须重新进入”。
\end{statementbox}

\section{成熟智能的工作记忆原则：选择性显式认知}

经过前面的分析，我们可以回到一个看似反直觉但实际上贯穿整个 AgentOS 理论的结论：\textbf{智能越成熟，越不应该把更多过程持续显式化。}

一个初学者面对陌生任务时，大量结构需要进入 h：
\[
NovelTask
\rightarrow
ExplicitRepresentation
\rightarrow
ExplicitPlanning
\rightarrow
ExplicitControl.
\]
随着经验累积，一部分结构逐渐形成：
\[
E,
\]
随后稳定为：
\[
\Theta,
\]
甚至形成更低成本的技能或外部程序。于是：
\[
ExplicitProcessing
\rightarrow
StableStructure.
\]

因此成熟过程通常表现为：
\[
\frac{C_h(\text{same task})}
{Performance}
\downarrow.
\]

这意味着真正的通用智能不能被理解为：
\begin{statementbox}
永远依靠更大的中央显式推理。
\end{statementbox}
更接近：
\begin{statementbox}
只在必要位置使用中央显式推理。
\end{statementbox}

我们可以把理想的认知资源配置写成一个优化问题：
\[
\min_{\pi_h}
\quad
\mathbb E
\left[
L_{\mathrm{task}}
+
\lambda_hC_h
+
\lambda_tT_{\mathrm{latency}}
+
\lambda_rR_{\mathrm{risk}}
\right],
\]
其中 $\pi_h$ 表示“什么结构何时进入工作记忆”的策略。

约束包括：
\[
C_h(t)\le C_h^{\max},
\]
以及：
\[
P(Success)\geq p_{\min},
\]
\[
Risk\leq R_{\max}.
\]

这个形式说明：减少 h 占用并不是目的本身。如果一个重要异常需要显式处理，却为了节省容量而被压在低层，系统会失败。因此成熟智能追求的不是：
\[
MinimumExplicitCognition,
\]
而是：
\begin{statementbox}
MinimumNecessaryExplicitCognition.
\end{statementbox}

也就是说：
\begin{statementbox}
该显式时必须显式，不需要显式时必须局部闭合。
\end{statementbox}

这与我们整个多尺度理论中的一个更一般原则一致：
\begin{statementbox}
NormalDynamicsStayLocal.
\end{statementbox}
以及：
\begin{statementbox}
UnexpectedResidualEscalates.
\end{statementbox}

工作记忆因此可以看成整个智能系统的“异常、冲突、新颖性和跨尺度协调面”。熟练的行为沉降以后不应持续占据 h；长期知识不应持续占据 h；全部身体状态不应持续占据 h；全部工具状态也不应持续占据 h。真正需要 h 的，是：

当前发生的新问题；

局部闭环无法解决的异常；

多个目标之间的冲突；

跨时间尺度结构需要重新协调的时刻；

重要自我状态需要显式进入的时刻；

需要产生新计划和新结构的时刻。

由此可以定义：
\[
h(t)
\approx
\mathcal C_{\mathrm{novel}}
\cup
\mathcal C_{\mathrm{conflict}}
\cup
\mathcal C_{\mathrm{residual}}
\cup
\mathcal C_{\mathrm{global}}
\cup
\mathcal C_{\mathrm{self-relevant}}.
\]

这一观点还解释了为何工作记忆有限性并不只是智能的限制。恰恰因为显式容量有限，智能系统才被迫发展出：

记忆层级；

技能；

抽象；

工具；

外化；

边界；

自动化；

局部闭环。

换句话说：
\[
\boxed{
FiniteWorkingMemory
\Rightarrow
PressureForCognitiveArchitecture.
}
\]

如果工作记忆无限大、处理成本为零，那么系统理论上可以把所有细节永久维持在当前空间中，也就没有强烈动力形成这些结构。但现实智能必然有限，因此智能的高级组织形式正是在有限性约束下出现的。

从这一意义看，工作记忆有限性不是一个需要“修复”的缺陷，而是整个智能架构形成的重要结构压力。

本章由此得到一个成熟智能判据：
\[
\boxed{
MatureIntelligence
=
SelectiveExplicitCognition
}
\]
即：
\begin{statementbox}
成熟智能不是一直在“想更多”，而是越来越准确地知道什么值得进入当前显式认知，什么应该由过去、结构、身体、工具和环境继续承担。
\end{statementbox}

\section{本章统一形式化：认知总线动力学}

最后，我们把本章内容压缩成一个统一的数学框架，以便后续章节能够在同一符号系统上继续发展。

定义 Agent 在时刻 $t$ 的认知结构全集：
\[
\mathcal C_A(t)
=
\mathcal C_W
\cup
\mathcal C_E
\cup
\mathcal C_\Theta
\cup
\mathcal C_\Psi
\cup
\mathcal C_{Me}.
\]

定义每个结构 $C_i$ 的当前状态：
\[
z_i(t)
=
\left(
r_i,
g_i,
n_i,
\varepsilon_i,
s_i,
c_i,
b_i
\right),
\]
其中分别表示相关度、目标相关性、新颖度、残差、自我相关性、显式处理成本以及边界状态。

定义提升函数：
\[
\Gamma_i(t)
=
F_\Gamma(z_i(t)).
\]

定义工作记忆准入：
\[
a_i(t)
=
\begin{cases}
1,&\Gamma_i(t)\geq\theta_i,\\
0,&\Gamma_i(t)<\theta_i.
\end{cases}
\]

则当前显式结构集合为：
\[
V_h(t)
=
\{C_i\mid a_i(t)=1\}.
\]

但由于有限容量，还必须满足：
\[
\mathcal C_h
\left(
V_h,E_h
\right)
\leq
C_h^{\max}.
\]

因此真正的准入问题应写成约束优化：
\[
V_h^*(t)
=
\arg\max_{V\subseteq\mathcal C_A}
\sum_{C_i\in V}
U_i(t)
\]
subject to
\[
C_h(V)\leq C_h^{\max}.
\]

其中 $U_i$ 可以综合：

目标价值；

风险；

新颖性；

预测残差；

自我相关性；

潜在信息收益。

工作记忆状态演化：
\[
\dot h
=
F_h
\left(
h,
\Pi_W(W),
\Pi_E(E),
\Pi_\Theta(\Theta),
\Pi_\Psi(\Psi),
\Pi_{Me}(Me)
\right).
\]

认知输出分成：
\[
h
\rightarrow
\begin{cases}
Action,\\
E,\\
\Theta,\\
Me,\\
Paused,\\
Closed.
\end{cases}
\]

因此总线的完整动力过程可以表示为：
\[
\boxed{
\mathcal C_A
\xrightarrow{\text{Promotion}}
h(t)
\xrightarrow{\text{Integration}}
h'(t)
\xrightarrow{\text{Disposition}}
\{Action,E,\Theta,Me,Pause,Close\}.
}
\]

其中 Promotion 解决“什么值得进入当前认知”；Integration 解决“当前结构怎样形成新的关系”；Disposition 解决“处理后的结构应该去哪里”。

这三步事实上构成了认知总线最基本的动力骨架。

进一步，从三相记忆角度：
\[
h
\]
是当前状态面；

\[
E
\]
是状态面留下的历史轨迹；

\[
\Theta
\]
是大量轨迹经过压缩后形成的生成结构。

因此：
\[
\boxed{
State
\rightarrow
Trajectory
\rightarrow
Generator
}
\]
和：
\[
\boxed{
Generator
\rightarrow
State
}
\]
不断通过 $h$ 闭合。

从世界交互角度：
\[
\boxed{
World
\rightarrow
h
\rightarrow
Action
\rightarrow
World
}
\]
不断通过 $h$ 闭合。

从自我角度：
\[
\boxed{
\Psi
\rightarrow
h^{self}
\rightarrow
Experience
\rightarrow
\Psi'
}
\]
也通过 $h$ 建立当前显式入口。

从内外记忆角度：
\[
\boxed{
Internal
\leftrightarrow
h
\leftrightarrow
External
}
\]
同样以 h 为交换面。

于是我们可以给本章作出最终定义：

\begin{definition*}
工作记忆 $h(t)$ 是一个容量有限、结构稀疏、动态重构的显式认知总线。它把世界观测、情景轨迹、长期生成结构、自我约束和外部认知载体选择性投影到同一个当前关系面，并通过行动、记忆写入、结构沉降、外化、暂停与闭合不断释放这一有限认知空间。
\end{definition*}

这一理解意味着，AgentOS 的核心并不是“拥有一个非常大的工作记忆”，而是拥有一套能够让正确结构在正确时刻进入工作记忆、并在处理完成后迅速离开工作记忆的动力机制。

因此，本章最终得到四条可以直接作为后续理论基础的原则：

\begin{proposition*}
第一，存在于 Agent 中的结构远多于显式存在于 $h$ 中的结构。
\end{proposition*}

\begin{statementbox}
第二，进入 $h$ 表示获得当前跨尺度显式处理权，而不是获得永久存在权。
\end{statementbox}

\begin{statementbox}
第三，成熟能力应减少对 $h$ 的持续占用，而异常、新颖与跨尺度冲突应重新获得进入 $h$ 的资格。
\end{statementbox}

\begin{statementbox}
第四，$h$ 的真正功能不是储存世界，而是在有限容量中组织世界、过去、自我与行动之间当前最重要的结构关系。
\end{statementbox}

由此，工作记忆从三相记忆体系中的“当前状态”进一步获得了更严格的系统论地位：

\[
\boxed{
h(t)
=
\textbf{the finite explicit cognitive bus of a persistent intelligent agent}.
}
\]

而一个持续智能体真正的成熟，不表现为它在每一个时刻都维持越来越庞大的显式思想，而表现为它能够把越来越多稳定结构留在适合它们的时间尺度和承载位置，只把此刻确实需要跨尺度协调的少量结构提升到这一有限的认知表面。
